\documentclass[conference]{IEEEtran}

\IEEEoverridecommandlockouts

% ---------------------------------------------------------------- packages
\usepackage{cite}
\usepackage{amsmath,amssymb}
\usepackage{graphicx}
\usepackage{booktabs}
\usepackage{array}
\usepackage{multirow}
\usepackage{tabularx}
\usepackage{adjustbox}
\usepackage{xcolor}
\usepackage{algorithm}
\usepackage{algpseudocode}
\usepackage{listings}
\usepackage{tikz}
\usetikzlibrary{shapes.geometric,arrows.meta,positioning,fit,calc,backgrounds}
\usepackage{textcomp}
\usepackage[hidelinks]{hyperref}

\newcommand{\code}[1]{\texttt{#1}}
\providecommand{\textasciicircum}{\^{}}

% --------------------------------------------------------------- listings
\definecolor{codegreen}{rgb}{0,0.5,0}
\lstdefinelanguage{MedScript}{
    morekeywords={input,const,func,let,if,else,while,return,true,false,int,float,bool,string,phi,range},
    morekeywords=[2]{kg,lb,gram,mg,ml,bpm,mmHg,degC,degF,metre,hour},
    morekeywords=[3]{print,export,present,anonymize,lb_to_kg,degF_to_degC,g_to_mg,abs,min,max},
    keywordstyle=\color{blue}\bfseries,
    keywordstyle=[2]\color{orange!80!black},
    keywordstyle=[3]\color{purple},
    sensitive=true, morecomment=[l]{//}, morestring=[b]"
}
\lstset{
    language=C, basicstyle=\ttfamily\fontsize{6.5}{7.5}\selectfont,
    keywordstyle=\color{blue}, commentstyle=\color{codegreen}, stringstyle=\color{violet},
    numbers=none, frame=single, breaklines=true, breakindent=0pt, showstringspaces=false,
    xleftmargin=0.2em, xrightmargin=0.2em, columns=fullflexible, aboveskip=4pt, belowskip=2pt
}

\newcommand{\NProg}{10{,}000}
\newcommand{\NFeat}{101}
\newcommand{\NTrain}{8{,}000}
\newcommand{\NTest}{2{,}000}
\newcommand{\MeanLoc}{38.7}
\newcommand{\TotalLoc}{386{,}953}
\newcommand{\NDouble}{1{,}398}
\newcommand{\DoubleShare}{15.5}
\newcommand{\RuleFirstAcc}{91.6}
\newcommand{\RulePrioAcc}{93.2}
\newcommand{\RuleFirstDouble}{65.7}
\newcommand{\RFAcc}{99.15}
\newcommand{\MLPAcc}{99.35}
\newcommand{\HGBAcc}{99.40}
\newcommand{\LRAcc}{99.20}
\newcommand{\RFMacro}{0.9915}
\newcommand{\RFDouble}{96.4}
\newcommand{\RuleLogicRecall}{68.5}
\newcommand{\RFLogicRecall}{97.0}
\newcommand{\MsPerKloc}{162}
\newcommand{\ScaleRSq}{0.981}
\newcommand{\ScaleSlope}{1.17}
\newcommand{\OptMean}{8.9}
\newcommand{\OptMax}{17.8}
\newcommand{\OptPrograms}{300}
\newcommand{\OptMismatch}{0}
\newcommand{\AugMinutes}{34}
\newcommand{\HeldoutNN}{7.4}


\begin{document}

\title{Intelligent Compiler for Detecting Errors in\\
Medical Data Processing Programs Using GAN\thanks{Course project for Compiler Design Lab (CSE3024), School of Computer
Science and Engineering (SCOPE), VIT-AP University. Compiler design component: parsing and semantic analysis.
Machine learning component: error classification. Application domain: healthcare.}}

\author{\IEEEauthorblockN{Swarnim Bandekar}
\IEEEauthorblockA{\textit{School of Computer Science and Engineering (SCOPE)}\\
\textit{VIT-AP University}\\
Amaravati, Andhra Pradesh, India\\
Regd. No. 24BCB7157}
\and
\IEEEauthorblockN{Lithin CG}
\IEEEauthorblockA{\textit{School of Computer Science and Engineering (SCOPE)}\\
\textit{VIT-AP University}\\
Amaravati, Andhra Pradesh, India\\
Regd. No. 24BCB7016}}

\maketitle

\begin{abstract}
Programs that process medical data can contain mistakes that a general-purpose compiler never reports: a dose added to a
body weight, a temperature threshold typed ten times too large, a laboratory value that may be missing but is used
unguarded, or a patient name that reaches a log file. This paper describes a compiler front end built to catch such
errors, and a learned classifier that names the root-cause family of each defective program. The front end is written for
MedScript, a small domain-specific language with units of measure in its types, declared physiological ranges,
protected-health-information tagging and optional values. It comprises a hand-written lexer, a recursive-descent parser with
error recovery, and a semantic analyser that checks types, units, range plausibility, taint flow of patient identifiers,
optional values, definite assignment and division by zero; a three-address-code generator and optimiser complete the
pipeline. On a synthetic corpus of \NProg{} programs in ten classes, classifiers using compiler-derived features reach
\RFAcc\% accuracy, against \RuleFirstAcc\% for the family of the first diagnostic, and they catch faults the compiler cannot see,
such as an inverted fever rule. A conditional generative adversarial network (GAN) trained to augment rare error classes did not
help: Gaussian jitter and SMOTE performed better, and only about 4\% of the GAN's samples were recognisable as the requested class.
A GAN discriminator trained on valid programs did detect the compiler-invisible faults well (AUROC 0.98). The corpus is synthetic,
and the results should be read accordingly.
\end{abstract}

\begin{IEEEkeywords}
compiler design, semantic analysis, parsing, units of measure, error classification, generative adversarial network,
healthcare software
\end{IEEEkeywords}
